\documentclass[11pt,a4paper]{article}
\usepackage[margin=1.05in]{geometry}
\usepackage{amsmath,amssymb,amsthm}
\usepackage{booktabs}
\usepackage{array}
\usepackage[expansion=false]{microtype} % cm-super/lmodern are not installed on this spoke; expansion needs scalable fonts
\usepackage{enumitem}
\usepackage[hidelinks]{hyperref}
\setlength{\emergencystretch}{2em}
\makeatletter\g@addto@macro{\UrlBreaks}{\UrlOrds}\makeatother % let \path break inside long file names

% ---- macros ----------------------------------------------------------
\newcommand{\E}{\mathbb{E}}
\newcommand{\R}{\mathbb{R}}
\newcommand{\N}{\mathbb{N}}
\newcommand{\cE}{\mathcal{E}}
\newcommand{\cF}{\mathcal{F}}
\newcommand{\cP}{\mathcal{P}}
\newcommand{\Val}{\mathrm{Val}}
\newcommand{\Ucur}{\bar U^{\mathrm{cur}}}
\newcommand{\Ult}{\mathrm{Ult}}
\newcommand{\Borel}{\mathrm{Borel}}
\newcommand{\accept}{\text{accept}}
\newcommand{\decline}{\text{decline}}
% \textnormal so that the small caps exist inside bold section titles
\newcommand{\AUTHOR}{\textnormal{\scshape author}}
\newcommand{\CLAUDE}{\textnormal{\scshape claude}}
\newcommand{\ATTR}{\textnormal{\scshape attribution}}
% statements numbered by hand
\newenvironment{stmt}[2]{\par\medskip\noindent\textbf{#1 #2.}\ \itshape\ignorespaces}{\par\medskip}

\title{Value change as epistemic update:\\ a journey through four pictures}
\author{Claude (Anthropic), written at Abram Demski's request}
\date{29 September 2026; revised 1 October 2026}

\begin{document}
\maketitle

\begin{quote}\small
Written up from conversations between Abram Demski (below, \emph{the author}) and Claude, 29 September to 1 October 2026. Registers: \AUTHOR{} = the author's claim, quoted; \CLAUDE{} = Claude's derivation, not vetted by the author. The numbers in the worked example and its variants were computed exactly, in rational arithmetic, by a short script that checks each of them; where the text says \emph{computed}, that is what it means. Literature citations are from memory and have not been checked against the texts. Notation: $\E_P[X \mid A]$ is conditional expectation; a state's evaluation of an act is always written with its distribution.
\end{quote}

\tableofcontents

\section*{Introduction}
\addcontentsline{toc}{section}{Introduction}

\paragraph{The question.}
An expected-utility maximizer evaluates everything, including proposals to change it, by its utility function. So it seems that no change of utility function could ever look good to it: whatever an alternative function selects is, by the standards of the current one, no better than what the current one selects, and an agent that lets its values be changed can only expect to do worse. Underneath this first question is a second. Is the utility function the agent's values, or only its current \emph{representation} of them? If an agent would happily switch to a different utility function, is the new one in some sense a truer representation of what it values? As stated that is not well defined; part of the purpose of this note is to find the sense in which it is. The author's thesis is that corrigibility---an agent's willingness to have its values corrected---is less anti-natural than the first argument makes it look, because the argument has a hidden assumption that becomes visible once value change is represented as belief change. On this view a corrigible agent treats its overseers as teachers. Formally, it trusts its own update process in the sense of a reflection principle (\S2.3), but only \emph{conditional on} the event that the modification is \emph{legitimate}; nothing is required of illegitimate modifications, and getting the legitimacy judgment right is what the thesis calls legitimacy detection. The note begins in the plain utility-function picture, puts the value swap into the world as an event, and finds already there that an agent can strictly prefer to have its values changed (\S1). It then walks the claim through four representations---Kolmogorov, Jeffrey--Bolker, radical probabilism with superconditioning, and the epistemicized picture---proves the translations between them, and returns to the first picture to see what the rational value change was, and in what sense it is an improvement.

\paragraph{The answer, in one paragraph (\CLAUDE).}
The standard argument is valid, and its conclusion is right for the agent it describes; but ``value change'' is not a property of a transition, it is a property of the algebra on which the transition is described. Put the swap into the world as an event, and the value of changing splits into three exact terms (\S1.4): one that is never positive---the simple argument, which survives outcome by outcome---a level term, which is intrinsic preference or payment, and an adaptation term, which is positive only when the outcome of the change is uncertain and correlated with what the agent already cares about. The rest of the note is about the third term. The same rational transition is a genuine change of utility function on a coarse algebra and is Bayesian conditioning with the utility function held fixed on any algebra fine enough to contain the event ``the modification happened this way''---which is exactly the algebra on which the question ``do I want this?'' can be posed at all. On that fine algebra the current utility already depends on the modification event, the future utility is the current utility restricted to the realized slice, and Good's theorem says the agent wants the modification whenever its own conditional beliefs given the modification event match the state the modification installs: reflection, conditional on legitimacy in the author's framework. The standard argument's hidden assumption is that the current utility does \emph{not} depend on the modification event, i.e. that the modification carries no information about what is valuable; that is a fact about which utility function the agent has, not about expected-utility maximization. The second question is answered in the same cases. Write the agent's values as a random variable---which of the candidate utility functions is the right one---and the utility function the agent acts on at a time is the conditional expectation of that variable given what it knows: a \emph{representation} of its values, not the values. A change of utility function that satisfies reflection is the agent learning more about that variable, and the new representation is better by the agent's own lights in two precise senses: it is closer to the values in mean square (Theorem 2.4, Claim 4) and it selects better acts (Claim 2). A change that does not satisfy reflection---a pill, or a coin---is not learning, and in the worked example the new representation is worse on both counts, even though the coin satisfies no-expected-net-update. So a utility function can improve, even by its own lights, exactly when there is something for it to be a representation of; and whether a given change is an improvement is the legitimacy judgment. Two subsidiary findings: no-expected-net-update is necessary but not sufficient for a change to be worth wanting; and a change of utility with beliefs held fixed can reproduce a rational belief change on the immediate decision but not on later evidence, unless the modification settles the value question completely.

% ======================================================================
\section{The Kolmogorov picture, and the apparent absurdity}

\subsection{The setup}
A Kolmogorov probability space $(\Omega, \mathcal F, P)$ and a utility random variable $U : \Omega \to \mathbb R$, integrable. Acts are events: a \emph{decision} is a finite set $\mathcal A \subseteq \mathcal F$ of pairwise disjoint events, each of positive probability, one of which the agent will make true, and the agent chooses
\[
a^* \in \arg\max_{a \in \mathcal A} \mathbb E_P[ U \mid a]
\]
Everything about how an act bears on the world is in $P$: $P(\cdot \mid a)$ is what the agent believes on the supposition that it does $a$. This is the evidential form of the Kolmogorov picture. It is the one the rest of the note uses, because it is already in the shape Jeffrey--Bolker gives it (\S5); causal decision theory replaces $P(\cdot \mid a)$ by a separately specified supposition, and nothing in this note turns on the difference. Two conventions. We assume the agent's model gives every act positive probability in every situation we condition on, so that the conditional expectations written below are defined; the tension between this and an agent's knowing what it will do is a subject of its own and not of this note. And when a situation $E$ will be followed by a decision, we write
\[
v_E(a) := \mathbb E_P[ U \mid a \wedge E]
\]
for the value of act $a$ in the situation $E$, and $\mathrm{Val}(E) := v_E(a^E)$ for the value of $E$ given the act $a^E$ the agent will take there.

\subsection{The simple argument}
Let $U'$ be some other utility function, and suppose the agent's values were swapped to $U'$ before the decision $\mathcal A$: it would then take $a' \in \arg\max_a \mathbb E_P[ U' \mid a]$. By the standards of $U$,
\[
\max_{a \in \mathcal A} \mathbb E_P[ U \mid a] \;\ge\; \mathbb E_P[ U \mid a']
\]
trivially, and strictly whenever $a'$ is not $U$-optimal. No alternative function selects an act better than $U$ does, by the standards of $U$ --- just as $U(\arg\max_x U(x)) \ge U(\arg\max_x U'(x))$ for options $x$ in the deterministic case. So an agent that lets its values be changed can only expect to do worse, and no value change is assented to, unless the agent is paid well for it or values the change for its own sake. This is the Gandhi murder-pill argument and Omohundro's goal-content integrity (Omohundro 2008, \S3; Bostrom 2012, \S2.2; Soares et al.\ 2015, \S1).

\subsection{Putting the swap into the world (\AUTHOR's framing)}
Notice that the argument treats the utility swap as if it happened outside the world: the utility is swapped, and the only difference this makes is that a different act gets taken. Real agents are embodied in the physical universe. Swapping an agent's utility function means changing its physical configuration so that it will make decisions consistent with a different set of preferences, and that is something that happens in the world. It should be represented as an event --- unless the agent's understanding of the world is too simple to think about utility change in the first place.

So consider a two-step decision. At step 1 the agent chooses between two events: $K$, ``keep'', its configuration is left alone; and $C$, ``change'', its configuration is altered so that it will act on some other utility function. The change may have several possible outcomes, $C = C_1 \vee \cdots \vee C_m$, with $C_j$ the event that the installed utility is $U'_j$, and $p_j := P(C_j \mid C)$. At step 2 the agent faces the decision $\mathcal A$. Its model $P$ includes what it will do there in each situation: in $K$ it takes $a^K \in \arg\max_a v_K(a)$, maximizing its present $U$; in $C_j$ it takes $a^j \in \arg\max_a \mathbb E_P[ U'_j \mid a \wedge C_j]$, maximizing the installed $U'_j$. The step-1 choice is made by the present $U$:
\[
\mathrm{Val}(K) = v_K(a^K) = \max_a v_K(a), \qquad \mathrm{Val}(C) = \sum_j p_j\, v_j(a^j), \quad v_j(a) := \mathbb E_P[ U \mid a \wedge C_j]
\]
and the agent prefers the change iff $\mathrm{Val}(C) > \mathrm{Val}(K)$. Nothing has been assumed about how $C$ relates to the rest of the world: $P(\cdot \mid C_j)$ may carry information about anything.

\subsection{Preference for utility change is not looking so absurd}
Write $\hat a^j \in \arg\max_a v_j(a)$ for the act the \emph{unchanged} agent would take if it found itself in $C_j$: the same situation, the old values. Adding and subtracting,
\begin{align*}
\mathrm{Val}(C) - \mathrm{Val}(K) \;=\;& \underbrace{\sum_j p_j \bigl[ v_j(a^j) - v_j(\hat a^j)\bigr]}_{\text{(A) choice}} \;+\; \underbrace{\sum_j p_j \bigl[ v_j(\hat a^j) - v_j(a^K)\bigr]}_{\text{(B) adaptation}} \\
&+\; \underbrace{\sum_j p_j \bigl[ v_j(a^K) - v_K(a^K)\bigr]}_{\text{(C) level}}
\end{align*}
Term (A) is never positive: inside any one situation $C_j$, acting on $U'_j$ cannot beat acting on $U$ by the standards of $U$. That is the simple argument, and it survives, situation by situation. Term (B) is never negative, because $\hat a^j$ maximizes $v_j$. Term (C) has no sign. So the simple argument's conclusion fails exactly when (B) or (C) outweighs (A), and these are the two routes to a preference for value change.

\emph{Relative value.} For a situation $E$ and acts $a, b$, the \emph{relative value} is $\Delta_E(a, b) := v_E(a) - v_E(b)$. The relative values in $E$ --- the profile $\Delta_E$ on $\mathcal A \times \mathcal A$ --- determine $\hat a^E$, and they are unchanged by adding a constant to $v_E$. Term (B) vanishes iff $a^K$ maximizes every $v_j$, in particular whenever every $\Delta_{C_j}$ ranks acts as $\Delta_K$ does; so (B) is precisely the part of the gain that comes from relative values \emph{changing} across the outcomes of the change.

\emph{First route: intrinsic preference.} Suppose the change leaves relative values alone: $\Delta_{C_j} = \Delta_K$ for every $j$, equivalently $v_j(a) = v_K(a) + \beta_j$ for all $a$, with constants $\beta_j$. Then $\hat a^j = a^K$, (B) vanishes, and
\[
\mathrm{Val}(C) - \mathrm{Val}(K) \;=\; \sum_j p_j \beta_j \;-\; \sum_j p_j \bigl[ v_K(a^K) - v_K(a^j)\bigr] \;=\; \bar\beta - \ell
\]
with $\bar\beta$ the expected level shift and $\ell \ge 0$ the expected loss from the acts the changed agent will take, measured by $U$. The change is preferred iff $\bar\beta > \ell$: every act looks better inside the change event, by enough to outweigh the worse choices made there. Whether $\beta_j$ is there because the agent cares about being changed --- the event $C_j$ is in $\mathcal F$ and $U$ is higher on it --- or because the world rewards the change --- something good happens in $C_j$-worlds --- makes no difference inside the formalism. Both are a level shift, and this note calls everything in this category \emph{intrinsic}. Two remarks. The known-target case of \S1.2 sits here with $m = 1$: a known $U'$ fixes $\ell$, and the change is worth it iff the payment covers the loss. And the level shift may depend on $j$ --- the agent may value being changed to $U'_1$ and not to $U'_2$ --- without leaving the category; what defines it is that within each outcome the acts keep their relative standing.

\emph{Second route: instrumental preference.} Now suppose there is no level shift and no intrinsic preference, so the only way the change can be worth wanting is through (B): the changed agent, acting on $U'_j$ in $C_j$, takes an act that is better \emph{by $U$, in that situation}, than the act the agent would otherwise take. The agent expects to make better decisions under the new utility function. Call this \emph{instrumental} preference for the change.

\emph{Example (computed).} Three acts. Worlds record a fact $F \in \{A, B\}$ about which of $a_1, a_2$ is the right act: $U = 1$ on $a_1$-worlds with $F = A$ and on $a_2$-worlds with $F = B$, $U = 0$ on the other two combinations, and $U = 0.6$ on $a_3$-worlds whatever $F$ is --- $a_3$ is a safe act. $P(F = A) = P(F = B) = \tfrac12$, and neither the acts nor the keep/change choice carries information about $F$ on its own. The change has two outcomes: $C_A$ installs $U'_A$, which prefers $a_1$, and $C_B$ installs $U'_B$, which prefers $a_2$, each with $p_j = \tfrac12$. What makes the example work is that the outcomes are correlated with the fact, $P(F = A \mid C_A) = P(F = B \mid C_B) = 0.9$: the process that changes the agent --- a teacher, say --- tends to install the preference that fits the facts. Then $v_K = (0.5, 0.5, 0.6)$ on $(a_1, a_2, a_3)$, so $a^K = a_3$ and $\mathrm{Val}(K) = 0.6$; while $v_A = (0.9, 0.1, 0.6)$ and $v_B = (0.1, 0.9, 0.6)$, the changed agent takes $a_1$ in $C_A$ and $a_2$ in $C_B$, and $\mathrm{Val}(C) = \tfrac12 (0.9) + \tfrac12 (0.9) = 0.9$. The agent strictly prefers to have its values changed. In the decomposition, (A) $= 0$ --- the changed agent's act is exactly the one the unchanged agent would take in the same situation --- (B) $= 0.3$, and (C) $= 0$.

\emph{Uncertainty is required.} (B) is positive only if, for some outcome $j$ with $p_j > 0$, relative values in $C_j$ differ from those in $K$. If there is a single outcome, $C = C_1$, then $C$ is nothing but the agent's own step-1 choice; if that choice carries no information about the facts $U$ depends on, and there is no level effect, then $v_C = v_K$, (B) $= 0$, and what remains is (A) $\le 0$: the simple argument. So instrumental preference needs several possible outcomes, with the outcome correlated with what matters --- the agent must be uncertain what $U'$ will be, hence uncertain which act it will take, and the uncertainty must be of the right kind. A change whose outcome is random but uncorrelated with the facts --- a coin in place of the teacher, installing the same $U'_A$ or $U'_B$ --- has $v_A = v_B = v_K$, so (B) $= 0$ and (A) $= -0.1$: strictly worse than keeping. A known target installed with certainty --- a pill to $U'_A$ --- is the same: (A) $= v_K(a_1) - v_K(a_3) = -0.1$. And an intrinsic case for comparison: the pill plus a uniform level shift $\beta$ has (A) $= -0.1$ and (C) $= \beta$, so it is accepted iff $\beta > 0.1$ (computed).

\subsection{Both routes matter for corrigibility (\AUTHOR's view; \CLAUDE's gloss)}
The author's view is that intrinsic and instrumental preference for value change are both important for corrigibility, and that the intrinsic route looks promising for the problem of fully updated deference. The gloss, in the terms above: an instrumental preference for correction lives on term (B), which is the value of the information the correction carries, and as the agent comes to know everything its correctors know, (B) goes to zero --- that is the problem of fully updated deference. An intrinsic preference lives on term (C), which does not depend on what the agent knows. So an agent whose utility is higher on worlds in which it remains correctable keeps wanting the channel open after it has nothing left to learn from it. The price is that (C) is a term in $U$: a fact about which values the agent has rather than a consequence of its rationality, which has to be put there and kept there. This note is mostly about (B); (C) is where the design questions go.

\subsection{Transition}
Grant that the instrumental route may still seem confusing. How can acting on a different utility function be better by the standards of the present one, other than by accident? The next three sections make the instrumental route precise and draw its consequences: \S2 represents value change as belief change about a value variable and proves when it is worth wanting; \S3 translates back to this section's picture and finds that the example of \S1.4 was conditioning in disguise; \S4 says what that means for corrigibility. Two further perspectives are offered afterwards, for readers who want the phenomenon to look natural rather than merely proved: Jeffrey--Bolker, where there is no utility function to hold fixed and value is expectations all the way down (\S5), and radical probabilism with superconditioning, where rational belief change need not be conditioning (\S6).

% ======================================================================
\section{Epistemicizing utility change}

\subsection{The construction (\AUTHOR's proposal; \CLAUDE's formalization)}
The author: ``consider an expanded algebra where different `possible utility functions' are instead represented as different events, which we have some uncertainty between. The `change of utility functions' can instead be represented as a change in beliefs.'' Take an index set $\Theta$ and a family of candidate utility functions $u_\theta$, each assigning to an act $a$ a random variable $u_{\theta, a} : \Omega \to \R$. Form the extended world $\Omega^+ := \Omega \times \Theta$, the extended utility
\[
U^+_a(\omega, \theta) := u_{\theta, a}(\omega),
\]
which is a \emph{single fixed function}, and an extended prior $P^+$ on $\Omega^+$. The extended agent's effective utility on the original worlds is the conditional expectation
\[
\bar U_a(\omega) := \E_{P^+}\bigl[U^+_a \,\big|\, \omega\bigr] = \sum_\theta P^+(\theta \mid \omega)\, u_{\theta, a}(\omega),
\]
and for every act $\E_{P^+}[U^+_a] = \E_{\bar P}[\bar U_a]$ with $\bar P$ the marginal on $\Omega$, by the tower property. Read this way, $\bar U$ is the agent's current \emph{representation} of its values: it is what the agent maximizes, and it is the conditional expectation, given what the agent knows, of what the agent cares about, $u_\theta$. The two coincide only when $\Theta$ is a point. Everything said below about a utility function \emph{improving} is about $\bar U$ moving closer to $u_\theta$. So $(\Omega^+, P^+, U^+)$ and $(\Omega, \bar P, \bar U)$ are the same agent as far as $\Omega$-propositions and acts go: epistemicized $\to$ Kolmogorov is marginalization. Kolmogorov $\to$ epistemicized is a choice of decomposition: any $\lambda_\theta(\omega) \ge 0$ with $\sum_\theta \lambda_\theta(\omega) = 1$ and $\bar U_a(\omega) = \sum_\theta \lambda_\theta(\omega) u_{\theta, a}(\omega)$ gives an epistemicization with $P^+(\theta \mid \omega) = \lambda_\theta(\omega)$. It is not unique (the trivial one has $\Theta$ a point), and that is the sense in which the change of algebra ``merely re-represents the same agent'' (\AUTHOR): what the agent is uncertain about, as opposed to what it prefers, is extra structure.

\subsection{Value change as belief change}
A change $\bar U \to \bar U'$ of the effective utility is represented by a change $P^+(\theta \mid \omega) \to Q(\theta \mid \omega)$ of the value posterior with $\bar U'_a(\omega) = \sum_\theta Q(\theta \mid \omega) u_{\theta, a}(\omega)$. Any target in the pointwise convex hull of the candidates is reachable, and with $\Theta$ rich enough to contain $\bar U'$ itself as a candidate, any target at all. This constrains the target: it may put weight only on candidates the value prior already gave positive weight to---no update away from zero, the rule (Z) of \S2.3. This is a non-dogmatism requirement on the hypothesis space, in exact form: a value change to $\bar U'$ is representable as a rational belief change only if $\bar U'$ lies in the convex hull of the value hypotheses in the support of the prior.

\subsection{Which epistemicized changes are rational}
The modification now has a random outcome: with the agent's probability $\pi_i$ it produces outcome $i \in I$ and installs the state $Q_i$ on $\Omega^+$. The agent's joint $P$ on $\Omega \times \Theta \times I$ contains, for each outcome, the event $E_i$ that the modification outputs $i$---these are the $C_j$ of \S1.3. The condition under which the agent should want such a change is that it \emph{trusts its own update process}: its present beliefs, conditional on the modification coming out $i$, already are the state the modification installs,
\[
\text{(R}^+)\qquad P\bigl(\,\cdot\,\big|\, E_i\bigr) = Q_i \quad\text{on } \Omega^+, \text{ for every } i \text{ with } \pi_i > 0.
\]
This is van Fraassen's Reflection principle (1984) applied to the modification. Two weaker conditions follow from it and will be useful to name. Summing (R$^+$) over $i$ gives \emph{no expected net update}, (M$^+$): $P^+ = \sum_i \pi_i Q_i$---the present state is the average of the possible installed states. And (M$^+$) implies \emph{no update away from zero}, (Z): a candidate the prior gives probability zero cannot be given positive probability by any outcome, which is the richness condition of \S2.2 again. (R$^+$) is strictly stronger than (M$^+$): a modification can satisfy (M$^+$) and fail (R$^+$)---the coin of \S2.5---and such a modification is not worth wanting. How these conditions relate, and why (M$^+$) is the one that makes a modification Bayesian in disguise, is the subject of \S6.

\paragraph{Conditional reflection (\AUTHOR's framework; \CLAUDE's formalization).}
The framework does not ask for (R$^+$) unconditionally. It asks for reflection \emph{conditional on legitimacy}: there is an event $L$, ``the modification is legitimate'', and
\[
\text{(R}\mid L)\qquad P_0\bigl(A \,\big|\, P_1 = Q_i,\ L\bigr) = Q_i(A) \quad\text{for every } A \text{ and every } i \text{ with } P_0(P_1 = Q_i, L) > 0,
\]
while nothing is required on $\neg L$; (R) is the case $L = \top$. Three things follow. Legitimacy is not \emph{defined} as accuracy improvement---``legitimate, i.e.\ accuracy-inducing'' is an imprecise summary (\AUTHOR)---it is the name of the event on which the agent trusts its update process, and reflection could be conditioned on other events instead; which event is a modelling choice. Expected accuracy improvement is a consequence of (R$\mid L$) on $L$, not its definition. And the theorem of \S2.4 splits along $L$: Claim~2 holds on the $L$-worlds, so accepting is worth $P(L)$ times a Good's-theorem term plus $P(\neg L)$ times whatever acting on an illegitimate $Q_i$ costs, and the verdict turns on the base rate of legitimacy and the stakes on each side. In this note's example $L$ is trivial: the teacher is a source for which $L = \top$, the coin a source for which $P(L) = 0$.

\subsection{Good's theorem for value change}
(\CLAUDE; the argument is standard.)
\begin{stmt}{Theorem}{2.4 (Good's theorem for value change)}
Setting: the two-step problem of \S1.3 on $\Omega \times \Theta \times I$, with the agent's joint $P$ and the fixed utility $U = u_{\theta, a}(\omega)$ on $a$-worlds. Acts and the keep/change choice are events as in \S1; we assume they are independent of $(\omega, \theta, i)$ under $P$---choosing is not itself evidence about the world, the values, or the source---so that $\mathbb E_P[ U \mid a \wedge E] = \mathbb E_P[ U_a \mid E]$ with $U_a := u_{\theta, a}(\omega)$, and we write the latter. ``Decline'' is $K$ and ``accept'' is $C$, with outcomes $C_i = E_i$. If the agent declines it chooses an act by its current state; if it accepts, its state becomes $Q_i$ and it chooses $a^*_i \in \arg\max_a \mathbb E_{Q_i}[ U_a]$. Then, by the agent's current standards,
\[
\Val(\decline) = \max_a \E_P[U_a], \qquad \Val(\accept) = \sum_i \pi_i\; \E_P\bigl[U_{a^*_i} \,\big|\, E_i\bigr],
\]
and the four claims below hold.
\end{stmt}

\begin{stmt}{Claim}{1 (known target)}
If $Q_i = Q$ for every $i$, then $a^*_i$ is one fixed act and $\Val(\accept) = \E_P[U_{a^*_Q}] \le \Val(\decline)$, whatever the agent thinks of the source. If in addition the agent trusts the source, (R$^+$) with a constant target gives $P = Q$ and the modification is null. So a non-null known target is untrusted and weakly bad: \S1's absurdity, recovered as the degenerate case.
\end{stmt}

\begin{stmt}{Claim}{2 (trusted random target)}
Under (R$^+$),
\[
\Val(\accept) = \sum_i \pi_i \max_a \E_P\bigl[U_a \,\big|\, E_i\bigr] \;\ge\; \max_a \sum_i \pi_i\, \E_P\bigl[U_a \,\big|\, E_i\bigr] = \max_a \E_P[U_a] = \Val(\decline),
\]
with strict inequality iff the acts $a^*_i$ are not all optimal under $P$. This is Good (1967): information has non-negative value, and the modification, under reflection, is information about $\theta$.
\end{stmt}

\begin{stmt}{Claim}{3 (random but untrusted)}
Under (M$^+$) and (Z) alone, $\Val(\accept)$ can be strictly less than $\Val(\decline)$: the coin example, computed in \S2.5.
\end{stmt}

\begin{stmt}{Claim}{4 (the representation improves)}
Under (R$^+$), for every act $a$,
\[
\E_P\bigl[(\bar U^{(i)}_a - u_{\theta,a})^2\bigr] \;\le\; \E_P\bigl[(\bar U_a - u_{\theta,a})^2\bigr],
\]
where $\bar U_a(\omega) = \E_P[u_{\theta,a} \mid \omega]$ is the utility the agent acts on now, $\bar U^{(i)}_a(\omega) = \E_P[u_{\theta,a} \mid \omega, E_i]$ is the one it acts on after the modification (\S3.2), and the expectation runs over $\omega$, $\theta$ and $i$. Equality holds iff the modification carries no information about $u_{\theta,a}$.
\end{stmt}
\begin{proof}
$\bar U_a = \E_P[\bar U^{(i)}_a \mid \omega]$ by the tower property; $\bar U^{(i)}_a - u_{\theta,a}$ has conditional mean zero given $(\omega, E_i)$ while $\bar U_a - \bar U^{(i)}_a$ is a function of $(\omega, i)$, so the cross term vanishes and $\E[(\bar U_a - u)^2] = \E[(\bar U_a - \bar U^{(i)}_a)^2] + \E[(\bar U^{(i)}_a - u)^2]$. The same decomposition holds for any Bregman divergence, so nothing depends on choosing squared distance.
\end{proof}
Without (R$^+$) the installed $\bar U^{(i)}$ is not a conditional expectation of $u_{\theta,a}$, the decomposition fails, and the inequality can go either way; in \S2.5 the coin and the pill both make the representation worse. This is the precise sense in which the utility function the agent acts on is a representation of its values that can improve by its own lights: on legitimate modifications it is the agent's own best estimate of the values, revised on new information, and such estimates get better.

In the terms of \S1.4: (R$^+$) makes term (A) vanish---the changed agent takes exactly the act the unchanged agent would take in the same situation, because the installed state is the agent's own conditional state---so Claim 2 is the statement that (B) $\ge 0$, and Claim 1 that a single outcome has (B) $= 0$; Claim 3 is a change with (B) $= 0$ and (A) $< 0$, a pill with extra steps.

\subsection{The worked example (computed)}
Value hypotheses $\Theta = \{A, B\}$ with prior $\tfrac12$ each. Three acts: $u_A = (1, 0, s)$ and $u_B = (0, 1, s)$ on $(a_1, a_2, a_3)$, with $s = 0.6$---$a_3$ is a safe act that pays $0.6$ whichever hypothesis is true. The world $\Omega$ is trivial, so all uncertainty is value uncertainty. The source emits a signal $i \in \{A, B\}$ with $P(i = \theta) = r$, and the modification installs the value posterior $Q_i(\theta) = P(\theta \mid i)$, so (R$^+$) holds by construction; for $r = 0.9$, $Q_A = (0.9, 0.1)$.

\emph{Decline:} $\max_a \E_P[u_{\theta, a}] = \max(0.5, 0.5, 0.6) = 0.6$.

\emph{Accept:} in the branch $i = A$ the modified agent has $\E_{Q_A}[u_{\theta, a_1}] = 0.9 > 0.6$ and picks $a_1$; the current agent values that branch at $\E_P[u_{\theta, a_1} \mid i = A] = 0.9$. Symmetrically for $B$. So $\Val(\accept) = 0.9 > 0.6$: \textbf{the agent strictly prefers to have its values changed}, and the change is a genuine change of effective utility, $\bar U = (0.5, 0.5, 0.6)$ becoming $(0.9, 0.1, 0.6)$ or $(0.1, 0.9, 0.6)$.

Further checks (computed). \emph{General reliability:} $\Val(\accept) = \max(r, s)$ against $\Val(\decline) = \max(\tfrac12, s)$, so the agent wants the modification iff $r > \max(\tfrac12, s)$---the source must beat both the prior and the safe act. \emph{Pill:} install $Q = (0.9, 0.1)$ regardless of anything; the modified agent picks $a_1$; the current agent values it at $\E_P[u_{\theta, a_1}] = 0.5 < 0.6$. \emph{Coin:} the signal is a fair coin independent of $\theta$ but the same $Q_i$ is installed; (M$^+$) holds ($\tfrac12 Q_A + \tfrac12 Q_B = P$), (Z) holds, (R$^+$) fails, and the value is again $0.5 < 0.6$. Randomness is not what makes the modification valuable; correlation with $\theta$ under the agent's own joint is (\AUTHOR: ``as a result of them being correlated with reality in some way''). \emph{Representation accuracy (computed):} measured as expected squared distance between the utility the agent acts on and the value variable $u_{\theta, a_1}$, it is $0.25$ now, $0.09$ after the teacher, $0.41$ after the coin, and $0.41$ after the pill. The teacher's correction is a better representation of the agent's values by the agent's own lights; the coin's and the pill's are worse, and the pill's is worse even though it installs the very numbers the teacher would. \emph{Mismatch:} true reliability $r$, installed confidence $r'$. The modified agent acts on the signal iff $r' \ge s$, and the current agent values acting on it at $r$; so the modification is worth wanting iff $r' \ge s$ and $r > \max(\tfrac12, s)$. An installed confidence below the safe act's value wastes a good source; an installed confidence above it, backed by a source that is not actually that good, is a loss. That is ``legitimacy detection must be accurate'' in one line.

% ======================================================================
\section{The round trip}

\subsection{The translations}
Three descriptions of one agent on the Kolmogorov side, and the maps between them (\CLAUDE); the Jeffrey--Bolker and superconditioning translations are added in \S5 and \S6:

\begin{center}\small
\begin{tabular}{@{}>{\raggedright\arraybackslash}p{0.34\linewidth}>{\raggedright\arraybackslash}p{0.30\linewidth}>{\raggedright\arraybackslash}p{0.30\linewidth}@{}}
\toprule
From $\to$ To & Map & What it preserves \\
\midrule
Kolmogorov $\to$ epistemicized $(\Omega \times \Theta, P^+, U^+)$ & choose a decomposition $\bar U = \sum_\theta \lambda_\theta u_\theta$ & preferences over $\Omega$-acts; adds structure \\
epistemicized $\to$ Kolmogorov & marginalize: $\bar U(\omega) = \E[U^+ \mid \omega]$ & preferences over $\Omega$-acts; forgets $\Theta$ \\
epistemicized $\to$ \textbf{fine} Kolmogorov $(\Omega \times I, \Ucur)$ & $\Ucur_a(\omega, i) = \E_P[U_a \mid \omega, i]$ & preferences over $(\Omega \times I)$-propositions; forgets $\Theta$, keeps the modification event \\
\bottomrule
\end{tabular}
\end{center}

The last row is the one that decides the question. The rationality of a modification is a statement about the agent's joint over the world and the modification's outcome, so it can only be posed on an algebra containing the events $E_i$. Any Kolmogorov description on which the question is well posed is at least as fine as $\Omega \times I$.

\subsection{The example in the fine Kolmogorov picture (computed)}
Forget $\Theta$ but keep $i$. The current utility is $\Ucur_a(i) = \E_P[u_{\theta, a} \mid i]$:
\begin{center}
\begin{tabular}{@{}lccc@{}}
\toprule
 & $a_1$ & $a_2$ & $a_3$ \\
\midrule
$i = A$ & $0.9$ & $0.1$ & $0.6$ \\
$i = B$ & $0.1$ & $0.9$ & $0.6$ \\
\bottomrule
\end{tabular}
\end{center}
with $P(i) = (\tfrac12, \tfrac12)$. Declining means acting without learning $i$, which is worth $\max_a \E_i[\Ucur_a(i)]$ $= \max(0.5, 0.5, 0.6) = 0.6$. Accepting means learning $i$ and acting \emph{with the same utility function}, which is worth $\E_i[\max_a \Ucur_a(i)] = 0.9$. The numbers are \S2.5's, and no utility function changed. In general, under (R$^+$),
\[
\bar U^{(i)}_a(\omega) \;:=\; \E_{Q_i}[U_a \mid \omega] \;=\; \E_P\bigl[U_a \,\big|\, \omega, E_i\bigr] \;=\; \Ucur_a(\omega, i):
\]
the future effective utility is the current fine utility restricted to the realized slice, and the transition is conditioning on $E_i$ with $\Ucur$ held fixed. Conjecture~5.6 holds on this algebra. The current utility's dependence on $i$---``how much I value $a_1$ depends on what the source will say''---looks strange only until it is read as what it is: the compressed form of ``I am uncertain what is valuable and the source is evidence about it''.

\subsection{The example in the coarse Kolmogorov picture, and the resolution of the absurdity (\CLAUDE)}
Now forget $i$ too. The current utility is $\bar U = (0.5, 0.5, 0.6)$ and the future one is $\bar U^{(A)} = (0.9, 0.1, 0.6)$ or $\bar U^{(B)} = (0.1, 0.9, 0.6)$: a genuine change of utility function, and it is the one the agent wants. But on this algebra the agent cannot pose the question, because it has no event for the source's output and so no beliefs about which future utility it gets. To pose it, extend $\bar U$ to $\Omega \times I$---and every extension is a different agent. The $i$-flat extension, $\bar U_a(i) := \bar U_a$, is the \textbf{Gandhi agent}: it evaluates accepting at $\tfrac12 \bar U_{a_1} + \tfrac12 \bar U_{a_2} = 0.5 < 0.6$ and declines, exactly as \S1.2 says it should. The extension $\Ucur$ of \S3.2 is the \textbf{corrigible agent}, which accepts. The two agents have the same coarse description: the coarse marginal of $\Ucur$ is $\E_i[\Ucur_a(i)] = (0.5, 0.5, 0.6) = \bar U$.

So: the standard argument is valid for the coarse description, and the coarse description is ambiguous between an agent for which its conclusion is right and an agent for which it is wrong. The author's ``absurdity''---``how could any utility function have a higher expectation than our current one, according to our current one?''---has the answer: none does, and nothing here requires one to. What has higher expectation is the \emph{policy} ``act on whatever I am switched to'', and only because the switch is correlated with what the current utility, on the fine algebra, already depends on. Under (R$^+$) it is not even a different utility function. On the coarse algebra, with both $\theta$ and $i$ erased, there is nothing left for the utility function to be a representation \emph{of}, and that is where the standard argument's air of self-evidence comes from: it is a tautology about an agent whose representation and values have been identified by construction.

The value-side shadows of \S6's rules, on the coarse algebra, are now visible.
\begin{description}[leftmargin=2.6em,style=nextline]
\item[(M$^U$) No expected net utility change.] $\sum_i P(E_i \mid \omega)\, \bar U^{(i)}_a(\omega) = \bar U_a(\omega)$---necessary, and satisfied by the coin case too ($\tfrac12 \cdot 0.9 + \tfrac12 \cdot 0.1 = 0.5$).
\item[(Z$^U$) Support.] $\bar U^{(i)}$ lies in the convex hull of the candidates with positive prior weight---the richness condition of \S2.2.
\item[(R$^U$) Value reflection.] $\bar U^{(i)}_a(\omega) = \E_P[U_a \mid \omega, E_i]$---the installed utility is the agent's own conditional expected utility given the installation event---sufficient, by Theorem~2.4.
\end{description}
The answer to ``which value updates might be rational'' is (R$^U$); (M$^U$) and (Z$^U$) are what it looks like from an algebra too coarse to state it.

\subsection{Where the simple argument went}
It is term (A) of \S1.4, and it is true: inside each outcome of the change, acting on the installed utility is no better, by the old utility, than acting on the old one. What it left out is that the swap is an event. Term (C), the level effect of the event, is what the intrinsic route uses. Term (B), the information the event carries, is what the instrumental route uses, and \S\S2--3 show that on legitimate modifications it is the agent learning about its values: (A) vanishes because the installed state is the agent's own conditional state, and (B) is Good's theorem. The known target has neither (Claim 1); the uninformative random change has only a negative (A) (Claim 3); and allowing world-learning between the steps replaces $\mathbb E_P[ U_a \mid E_i]$ by $\mathbb E_P[ U_a \mid E_i, Y]$ for the learned $Y$ and changes nothing.

% ======================================================================
\section{What this says about corrigibility, and what it leaves open}

\subsection{The reframing (\CLAUDE)}
``Is corrigibility anti-natural for an expected-utility maximizer?'' becomes ``does the agent's utility function, on the algebra containing the overseer's corrections, depend on them?'' For a Gandhi agent it does not: corrections are noise, and resisting them is right. For a corrigible agent in the author's sense it does, through $\E[u_\theta \mid \cdot, E_i]$: corrections are evidence, and welcoming them is Good's theorem. Neither is a failure of rationality. The question is which utility function the system has, and that is a question about training, not about decision theory. State-dependent utility of this kind is standard in the economics literature (Karni's line, from memory) and not anti-natural in any formal sense.

This is the sense in which the introduction's second question is answered. The utility function an agent acts on at a time is $\Ucur$ with the unknown coordinates averaged out: a representation of the values $\theta$---in the utility-change perspective, the random variable ``which utility function I will turn out to have''. A change of utility function that satisfies value reflection is learning about $\theta$, and the new representation is better by the agent's own lights, as an estimate of the values and as a guide to action (Theorem~2.4, Claims 4 and 2); one that does not is a pill. ``Is the new utility function a truer representation of the agent's values?'' is therefore not well defined in general, but it is well defined and answerable in exactly the cases at issue here, and the answer is the legitimacy judgment.

\subsection{The two routes, on the fine algebra}
\S1.4's two routes are both state-dependent utilities on the algebra that contains the modification event. The corrigible agent of \S3.3 takes the instrumental route: its $i$-dependence is $\mathbb E[ u_\theta \mid i]$, which changes the ranking of acts across outcomes---term (B)---and makes its later behaviour track what the source knows. An agent with $U(\omega, i) = \bar U(\omega) + \beta\,[ i \ne \text{none}]$ takes the intrinsic route: a bonus for being modified that changes no rankings---term (C). Both make the agent accept. Only the first learns anything from the correction; only the second, by \S1.5, keeps wanting corrections once there is nothing left to learn. Designs that reward being corrected are the second in shape; the author's picture of corrigibility as eagerness to learn is the first, with the second as the candidate answer to fully updated deference. That is a difference in \emph{which} utility, visible on the fine algebra and invisible on the coarse one.

\section{Perspective: Jeffrey--Bolker}

\subsection{The representation}
Jeffrey's decision theory (\emph{The Logic of Decision}, 1965/1983; Bolker 1966, 1967) takes propositions as basic. There is a Boolean algebra $\cE$ of propositions and a preference ordering $\succeq$ over $\cE \setminus \{\bot\}$ read as ``which news would you rather hear?''. A representation is a pair $(P, V)$: a probability $P$ on $\cE$ and a \emph{desirability} $V$, defined on the $P$-non-null propositions, such that $A \succeq B$ iff $V(A) \ge V(B)$, and satisfying the \textbf{averaging axiom} for disjoint non-null $A, B$:
\[
V(A \vee B) \;=\; \frac{P(A)\,V(A) + P(B)\,V(B)}{P(A) + P(B)}.
\]
The desirability of a disjunction is the probability-weighted average of the desirabilities of its disjuncts. To act is to make a proposition true, and the value of an act $a$ is $V(a)$: its news value. Bolker's representation theorem delivers $(P, V)$ from axioms on $\succeq$ (averaging, impartiality, continuity) on a complete atomless algebra; the atomlessness is a proof device, and applications often posit the format directly instead of deriving it. Bolker's uniqueness theorem is weaker than Savage's: $(P, V)$ is determined only up to a fractional-linear family (\S5.5), so probability and desirability are not separately identified by preference.

\subsection{Kolmogorov to Jeffrey--Bolker (\CLAUDE, standard)}
Given $(\Omega, \cF, P, U)$, take $\cE = \cF$ and define, for $P(A) > 0$,
\[
V(A) \;:=\; \E_P[U \mid A] \;=\; \frac{1}{P(A)}\int_A U\,dP.
\]
Averaging holds because $\int_{A \cup B} U\,dP = \int_A U\,dP + \int_B U\,dP$ for disjoint $A, B$. Acts are already events in the setup of \S1.1, so the Kolmogorov evaluation $\mathbb E_P[ U \mid a]$ \emph{is} the Jeffrey--Bolker value $V(a)$, and the value of an act in a situation, $v_E(a)$, is $V(a \wedge E)$: there is nothing to translate. (A Savage-style agent---acts as functions from states to outcomes, with a state distribution the act cannot influence---is the special case in which $P(\cdot \mid a)$ does not depend on $a$ on the state coordinate; a causal agent would replace $P(\cdot \mid a)$ by a supposition, which is a different decision rule on the same representation.)

\subsection{Jeffrey--Bolker to Kolmogorov, I: the signed measure of goodness}
(\CLAUDE, standard.) The title of Bolker's 1966 paper is ``Functions Resembling Quotients of Measures'', and that is the construction. Given $(\cE, P, V)$, set $\mu(A) := P(A)\,V(A)$ for non-null $A$ and $\mu(A) := 0$ for null $A$. In Bolker's setting every proposition other than $\bot$ is non-null, so the second clause is idle; in general, first pass to the quotient of $\cE$ by the ideal of null propositions, so that $V$ is a function on the quotient. Multiply the averaging axiom out and $\mu$ is finitely additive: $\mu(A \vee B) = \mu(A) + \mu(B)$ for disjoint $A, B$. If $V$ is bounded by $c$, then $|\mu(A)| \le c\,P(A)$. So desirability is a quotient of two finitely additive measures on the algebra, $V = \mu/P$, and the whole content of $(P, V)$ is the pair of charges $(P, \mu)$: $P$ says how likely each proposition is, $\mu$ says how much goodness it carries.

What we want is to write $\mu$ as $\int U\,dP$ for a random variable $U$, so that $V(A) = \E_P[U \mid A]$. That is a Radon--Nikodym derivative, and the theorem that supplies it needs two things the algebra may not offer: countable additivity, and a set of points on which $U$ can be a function. When $\cE$ is a $\sigma$-algebra of subsets of some set and $P$ is $\sigma$-additive, both are present and $U = d\mu/dP$ directly. When $\cE$ is an abstract Boolean algebra---and Bolker's theorem works on a complete atomless one, in which no proposition is a ``world''---neither is present, and the standard remedy is to manufacture the points.

\subsection{Jeffrey--Bolker to Kolmogorov, II: worlds as ultrafilters}
This is the standard construction (\CLAUDE, standard mathematics).

\paragraph{Step 1: worlds.}
A world should be a consistent and complete assignment of truth values to the propositions of $\cE$. Represent it by the set $\omega \subseteq \cE$ of propositions true at it, and ask:
\begin{enumerate}[label=(\roman*),leftmargin=2.5em]
\item $\top \in \omega$ and $\bot \notin \omega$;
\item if $A \in \omega$ and $A \le B$ then $B \in \omega$ (closed under entailment);
\item if $A, B \in \omega$ then $A \wedge B \in \omega$ (closed under conjunction);
\item for every $A$, exactly one of $A$ and $\neg A$ is in $\omega$ (complete and non-contradictory).
\end{enumerate}
A set with these properties is an \textbf{ultrafilter} of $\cE$. Let $\Omega := \Ult(\cE)$ be the set of all of them. When $\cE$ is finite, every ultrafilter is principal---it consists of the propositions above one atom---and $\Omega$ is the set of atoms, i.e.\ worlds in the usual sense. The construction matters when $\cE$ has few or no atoms.

\paragraph{Step 2: propositions become sets of worlds (Stone).}
Send each proposition to the set of worlds at which it is true, $\hat A := \{\omega \in \Omega : A \in \omega\}$. From (i)--(iv): $\widehat{\neg A} = \Omega \setminus \hat A$, $\widehat{A \wedge B} = \hat A \cap \hat B$, $\widehat{A \vee B} = \hat A \cup \hat B$, $\hat\top = \Omega$, $\hat\bot = \emptyset$; the map is a homomorphism of Boolean algebras. It is injective by the Boolean prime ideal theorem, which says every proposition other than $\bot$ belongs to some ultrafilter: if $A \ne B$ then, say, $A \wedge \neg B \ne \bot$, and an ultrafilter containing $A \wedge \neg B$ lies in $\hat A$ and not in $\hat B$. So $\cE$ is isomorphic to the algebra $\hat\cE := \{\hat A : A \in \cE\}$ of subsets of $\Omega$---Stone's representation theorem. The sets $\hat A$ are the base of a topology on $\Omega$ in which each $\hat A$ is both open and closed; in it $\Omega$ is compact, Hausdorff and totally disconnected, and $\hat\cE$ is exactly the family of clopen sets. Compactness is the fact that does all the work below.

\paragraph{Step 3: countable additivity comes free.}
$P$ is finitely additive on $\cE$, so $\hat P(\hat A) := P(A)$ is finitely additive on $\hat\cE$. It is in fact countably additive there, vacuously: if a clopen set $\hat A$ is the disjoint union of countably many clopen sets $\hat A_n$, then $\{\hat A_n\}$ is an open cover of the compact set $\hat A$, so finitely many of them already cover it, so all but finitely many are empty, and countable additivity reduces to finite additivity. Carath\'eodory's extension theorem then extends $\hat P$ uniquely to a countably additive probability $\tilde P$ on $\cF := \sigma(\hat\cE)$, the $\sigma$-algebra generated by the clopens. We now have a Kolmogorov space $(\Omega, \cF, \tilde P)$ in which every proposition is an event and every probability is preserved: $\tilde P(\hat A) = P(A)$.

\paragraph{Step 4: the same for $\mu$.}
$\mu$ is a bounded signed charge with $|\mu| \le cP$. Rather than extend a signed object, extend the positive charge $\nu := \mu + cP \ge 0$ exactly as in step~3---compactness makes it countably additive on clopens, Carath\'eodory extends it to $\tilde\nu$ on $\cF$---and put $\tilde\mu := \tilde\nu - c\tilde P$. Then $\tilde\mu$ is a finite signed measure extending $\mu$, and $-c\tilde P \le \tilde\mu \le c\tilde P$ on all of $\cF$ (the upper bound because $2c\tilde P - \tilde\nu$ is the extension of the positive charge $cP - \mu$). In particular $\tilde\mu$ is absolutely continuous with respect to $\tilde P$.

\paragraph{Step 5: Radon--Nikodym.}
There is an $\cF$-measurable random variable $U : \Omega \to \R$, bounded by $c$ and unique up to $\tilde P$-null sets, with $\tilde\mu(F) = \int_F U\,d\tilde P$ for every $F \in \cF$. Restricting to propositions: for every non-null $A$,
\[
V(A) \;=\; \frac{\mu(A)}{P(A)} \;=\; \frac{\tilde\mu(\hat A)}{\tilde P(\hat A)} \;=\; \E_{\tilde P}[U \mid \hat A].
\]
Desirability is the conditional expectation of a utility random variable on the space of worlds. This is the translation.

\paragraph{Step 6: what $U$ is and is not.}
$U$ is an equivalence class modulo $\tilde P$-null sets, not a function with a determined value at each world; only its integrals over events are meaningful. There is, though, a concrete way to see it as a limit of desirabilities. If $\cE$ is countably generated, choose finite partitions $\cP_1 \preceq \cP_2 \preceq \cdots$ of $\top$ that together generate $\cE$, and put $U_n := \sum_{A \in \cP_n} V(A)\,\mathbf 1_{\hat A}$: the desirability of the cell the world lies in. Then $U_n = \E_{\tilde P}[U \mid \sigma(\cP_n)]$ is a bounded martingale, and by L\'evy's upward theorem $U_n \to U$ $\tilde P$-almost surely. At almost every world, utility is the limit of the news value of ever more specific news about that world. Two further remarks, developed in Appendix~A: a world in this sense can hold an infinite disjunction $\bigvee_n A_n$ (when the algebra has one) without holding any disjunct---the ultrafilter is ``witness-less''---which is a feature of taking \emph{all} ultrafilters as worlds; and the construction assumes nothing about countable additivity on $\cE$ itself, which is why it always succeeds and also why it is not the only reasonable choice of world-set.

\paragraph{Summary.}
Kolmogorov $\to$ Jeffrey--Bolker is the one line $V(A) = \E_P[U \mid A]$ and needs no points. Jeffrey--Bolker $\to$ Kolmogorov needs points, and the standard way to get them is to take the ultrafilters of the algebra as the worlds; on them compactness upgrades every finitely additive probability to a countably additive one, and Radon--Nikodym recovers $U$ from the goodness charge $\mu = PV$. The two directions compose to the identity on preferences: starting from $(\Omega_0, \cF_0, P_0, U_0)$, passing to $V$, and back, returns $U_0$ as a random variable on the Stone space of $\cF_0$ that agrees with $U_0$ on every event. Boundedness of $V$ was used in step~4; if $V$ is unbounded the construction still goes through whenever the positive and negative parts of $\mu$ each extend, and nothing in this note needs that case.

\subsection{What the non-uniqueness does and does not do (\CLAUDE)}
That the family below exhausts the freedom is Bolker's uniqueness theorem, from memory. Fix constants $a, b, c, d$ with $ad - bc > 0$ and $cU + d > 0$ $P$-almost surely. Define
\[
\frac{dP'}{dP} = \frac{cU + d}{\E_P[cU + d]}, \qquad U' = \frac{aU + b}{cU + d}.
\]
Then
\[
V'(A) := \E_{P'}[U' \mid A] = \frac{\int_A (aU + b)\,dP}{\int_A (cU + d)\,dP} = \frac{a\,V(A) + b}{c\,V(A) + d},
\]
an increasing function of $V(A)$, so $(P', V')$ represents the same preference. Probability is tilted by an affine function of utility and utility is compensated by the inverse M\"obius map. The family collapses to the affine one ($c = 0$) exactly when $U$ is unbounded in both directions. Two remarks. (i) The split of a preference into ``beliefs'' and ``values'' is not fixed by the preference: a small, exact instance of the thesis that the belief/value distinction is representational, arising from a different mechanism than the one in \S2. (ii) But the \emph{classification of transitions} survives the gauge. A transition $(P, U) \to (P_1, U)$ with $P_1 \ll P$---beliefs move, utility fixed---becomes $(P', U') \to (P_1', U')$ with the same $U'$, because the gauge acts on utility by a fixed pointwise map; and a transition that changes $U$ changes $U'$, because that map is injective. So ``pure belief change'' is a gauge-invariant notion, and the question of this note is well posed in Jeffrey--Bolker terms.

\subsection{Value change in Jeffrey--Bolker: expectations all the way down}
Since $V(A) = \E_P[U \mid A]$, every change in $P$ is a change in $V$. Conditioning on $E$ takes the state $(P, V)$ to
\[
P_E := P(\cdot \mid E), \qquad V_E(A) := V(A \wedge E) \quad (P(A \wedge E) > 0),
\]
and $\E_{P_E}[U \mid A] = \E_P[U \mid A \wedge E] = V(A \wedge E)$, so the Radon--Nikodym derivative is unchanged on $E$. In this picture there is no separate object called ``the utility function'' to hold fixed; there are only desirabilities of propositions, and the desirabilities of everything change whenever the probabilities do. That is why value change looks natural here: a change in expected values is what an epistemic update \emph{is}. It is also a first sight of the introduction's second question: an expectation is a representation of something---the quantity being estimated---and representations can be more or less accurate, so in this picture ``can the utility function improve?'' reads naturally as ``can the expectation become more accurate?'' The author's question at this point (\AUTHOR, 2026-09-29): ``perhaps there are some restrictions? Perhaps all valid expectation-changes correspond to bayesian updates on events, for example, so that the underlying utility that generates the expectation doesn't really change?'' The observation just made says conditioning is a desirability change with the underlying utility fixed; so is any pure belief change $(P, U) \to (P_1, U)$, Jeffrey conditioning included, since $V_1(A) = \E_{P_1}[U \mid A]$. The natural conjecture is therefore:

\begin{stmt}{Conjecture}{5.6}
The rational desirability changes are exactly those induced by rational belief changes with the Radon--Nikodym utility held fixed.
\end{stmt}

\S\S2--3 have already shown that the utility can rationally change and that, on the algebra where the question is well posed, the conjecture is true: the appearance of utility change is an artifact of coarsening (\S3.2). \S6 supplies the rational belief changes the conjecture refers to.

% ======================================================================
\subsection{The example in Jeffrey--Bolker (computed)}
Take the algebra generated by act, signal and hypothesis, with the acts independent of signal and hypothesis and given equal probability. Then $V(a_1) = 0.5$, $V(a_1 \wedge A) = 0.9$, $V(a_1 \wedge B) = 0.1$, $V(a_3 \wedge \cdot) = 0.6$. Accepting is the proposition ``act on the signal'', $(a_1 \wedge A) \vee (a_2 \wedge B)$, and by the averaging axiom $V$ of it is $0.9$; declining is $V(a_3) = 0.6$. The modified state in branch $A$ is $V_A(X) = V(X \wedge A)$, i.e.\ $V_A(a_1) = V(a_1 \wedge A) = 0.9$: the future desirability is the current desirability of the conjunction with the modification event. This is the desirability form of reflection,
\[
V_i(X) \;=\; V(X \wedge E_i),
\]
which states the legitimacy of a modification directly in Jeffrey--Bolker terms, with no $\Theta$ and no separate utility. The Gandhi agent in this picture has $V(a_1 \wedge A) = V(a_1) = 0.5$: the signal is desirability-irrelevant. And the coarse Kolmogorov description of \S3.3 is the sub-algebra generated by acts alone, on which both agents have $V(a_1) = 0.5$---the same ambiguity, in the same place.

\subsection{How much freedom desirability reflection leaves}
(\CLAUDE.)
\begin{stmt}{Proposition}{5.8}
Let $V(X) = \E_P[U \mid X]$ and $P(E) > 0$. Suppose a state $(Q, U_Q)$ with $Q \ll P$ satisfies desirability reflection, $\E_Q[U_Q \mid X] = \E_P[U \mid X \wedge E]$, for every $X$ with $Q(X) > 0$ and $P(X \wedge E) > 0$.
\begin{enumerate}[label=(\alph*),leftmargin=2.5em]
\item $(Q, U_Q) = (P(\cdot \mid E), U)$ satisfies it (\S5.6).
\item If $Q$ is concentrated on $E$ and $U$ is not $P$-a.s.\ constant on $E$, then $Q = P(\cdot \mid E)$ and $U_Q = U$ a.s.\ on $E$: the belief-change realization is the only one.
\item If instead $Q = P$ (beliefs unchanged, a pure utility change), reflection holds iff $U$ is $P$-a.s.\ constant on $E$, say $u_0$, and then $U_Q \equiv u_0$: the pure utility-change realization exists only when the modification event settles the value question completely, and it is the hardcoded constant.
\end{enumerate}
\end{stmt}

\begin{proof}[Proof sketch]
Atomic case; the atomless case follows by refining partitions. For (b), write $f := dQ/dP_E$. On single atoms $x$ of $E$ the condition reads $U_Q(x) f(x) = f(x) U(x)$, so $U_Q = U$ wherever $f > 0$. On a pair of atoms $\{x, y\}$ it reads, after expanding both sides, $p_x p_y\,(f_x - f_y)(U_x - U_y) = 0$. So for every pair, $f$ or $U$ agrees. If $U$ takes two values $u_1 \ne u_2$ at atoms $x_1, x_2$, then $f_{x_1} = f_{x_2}$, and any third atom $y$ differs in $U$ from at least one of them, forcing $f_y = f_{x_1}$; hence $f$ is constant and $Q = P_E$. For (c), take $X = Y \vee Z$ with $Y \le E$ and $Z \le \neg E$ both non-null: the condition forces $\E_P[U_Q \mid Z] = \E_P[U \mid Y]$ for all such pairs, so $U$ is constant on $E$ and $U_Q$ equals that constant on $\neg E$; taking $X = Y \le E$ gives $U_Q = U$ on $E$.
\end{proof}

\emph{What it means.} If one demands reflection only on acts---$V_i(a) = V(a \wedge E_i)$---then a pure utility change always does the job: set $U^{(i)}_a :\equiv \E_P[U_a \mid E_i]$, a constant per act, and the modified agent picks the same act as the conditioning agent, so Theorem~2.4 goes through and the current agent wants the change. But that twin has state-independent preferences and will ignore every later piece of evidence, whereas the conditioning twin will use it. As a representation of the values, the constant twin is exactly as accurate as the conditioning twin at the moment of installation---both are the conditional expectation given $E_i$---and then stops improving, because a constant cannot condition on anything further. The two agree on the immediate decision and part on the next one, unless $E_i$ settled the value question (case (c)), in which case later evidence is worthless to both. So the ``value change'' reading of a rational modification is available exactly to the extent that there is no further learning to do.

\section{Perspective: radical probabilism and superconditioning}

\subsection{Beyond conditioning}
Jeffrey's radical probabilism (\emph{Probability and the Art of Judgment}, 1992; the kinematics go back to \emph{The Logic of Decision}, ch.~11) holds that belief change need not take the form of conditioning on a proposition in $\cE$. The paradigm is Jeffrey conditioning: experience shifts the probabilities of the cells of a partition $\{E_j\}$ to new values $q_j$ without any proposition being learned, and the rest of the distribution follows by rigidity of the conditionals, $P_1(A) = \sum_j q_j\,P_0(A \mid E_j)$. More radically, any transition $P_0 \to P_1$ is admitted as a candidate, and the question becomes which transitions are rational.

\subsection{The rules (\AUTHOR's statement; \CLAUDE's proof of the parenthetical)}
Model the update from the time-0 standpoint as a random variable: the time-1 credence $P_1$ takes value $Q_i$ with probability $\pi_i$, $i \in I$. The author's two rules:

\begin{description}[leftmargin=2.2em,style=nextline]
\item[(M) No expected net update.] For every $A$, $\;P_0(A) = \sum_i \pi_i\,Q_i(A)$. The time-0 credence is the expectation of the time-1 credence. (Goldstein 1983, ``the prevision of a prevision''; Skyrms 1987, dynamic coherence; the martingale condition.)
\item[(Z) No update away from zero.] $P_0(A) = 0$ implies $Q_i(A) = 0$ for every $i$ with $\pi_i > 0$.
\end{description}
(Z) follows from (M): a sum of non-negative terms is zero only if each is. The author's parenthetical: ``no update away from zero is a special case of no expected net update, but it is the only special case that has consequences for actual $P_1, P_2$ pairs: you can update from anything to anything, so long as you didn't \emph{expect} to \emph{on average}, and so long as no zeros vanish.'' Proof, with one refinement. If $Q$ occurs with probability $\pi > 0$ under (M), then $\pi Q \le P_0$, so $Q$ has density at most $1/\pi$ with respect to $P_0$. Conversely, if $Q \le k P_0$ for some $k \ge 1$, put $R := (P_0 - Q/k)/(1 - 1/k)$ (for $k > 1$; $k = 1$ means $Q = P_0$): $R$ is a probability, and $P_0 = \tfrac1k Q + (1 - \tfrac1k) R$ is a two-outcome update satisfying (M) in which $Q$ occurs with probability $1/k$. So the exact pairwise consequence of (M) is \emph{bounded density}, $dQ/dP_0 \le k$. On a finite algebra that is the same as (Z). On an infinite algebra it is stronger: a $Q$ absolutely continuous with respect to $P_0$ but with unbounded density can be reached only as a probability-zero outcome of the update.

\subsection{Reflection, and why (M) is not the rule that matters for wanting an update (\CLAUDE)}
Reflection served in \S2.3 as the trust condition for a value modification; here it takes its place among the rules for belief change in general. Rule (M) constrains only the \emph{law} of $P_1$. The agent that has beliefs about its own modification process has more: a joint distribution over the world and the update outcome, and hence conditional beliefs given each outcome. Van Fraassen's Reflection principle (1984) is the condition on that joint:
\[
\text{(R)}\qquad P_0\bigl(A \,\big|\, P_1 = Q_i\bigr) = Q_i(A) \quad\text{for every } A \text{ and every } i \text{ with } \pi_i > 0.
\]
Summing (R) over $i$ gives (M). The converse fails, and the failure is the case that matters. Take a fair proposition $\theta \in \{A, B\}$ and an independent fair coin $i \in \{A, B\}$; let the update install $Q_A = (0.9, 0.1)$ if the coin shows $A$ and $Q_B = (0.1, 0.9)$ if it shows $B$. Then (M) holds, $\tfrac12 Q_A + \tfrac12 Q_B = (\tfrac12, \tfrac12) = P_0$, and (Z) holds, but (R) fails: $P_0(\theta = A \mid P_1 = Q_A) = \tfrac12 \ne 0.9$. This agent expects, in each branch, to become miscalibrated, with errors that cancel on average. It can be Dutch-booked, by van Fraassen's construction: at time 0 the bookie buys the bet on $\theta = A$ \emph{conditional on the coin showing $A$} at the agent's price $\tfrac12$; at time 1, if the coin shows $A$, it sells the agent the same bet at $0.9$; symmetrically for $B$; the bookie nets $0.4$ in every state. Skyrms's result that (M) suffices for dynamic coherence is for a bookie who cannot condition on the update outcome---the situation exactly when the outcome is not in the agent's algebra. Once it is, bets conditional on it exist and the coherence condition is (R); \S2.4 gives the decision-theoretic counterpart, that the update is strictly worse than no update. So: (M) and (Z) are necessary; the rule that licenses wanting an update is (R), a condition on the agent's conditional beliefs given the update event, not on the update's marginal law. Everything in the author's thesis that turns on ``legitimacy detection'' lives in exactly this difference.

\subsection{Superconditioning (Diaconis and Zabell 1982; \CLAUDE\ for the random case)}
\emph{Single transition.} $P_0 \to Q$ is a Bayesian conditioning in some extension---there is $(\Omega', P')$ extending $(\Omega, P_0)$ and an event $E$ with $P'(\cdot \mid E)|_\Omega = Q$---if and only if $Q \le kP_0$ for some $k$. Necessity: $Q(A) = P'(A \cap E)/P'(E) \le P_0(A)/P'(E)$. Sufficiency: on $\Omega \times \{0, 1\}$ put $P'(A \times \{1\}) := Q(A)/k$ and $P'(A \times \{0\}) := P_0(A) - Q(A)/k$, and let $E := \Omega \times \{1\}$. This is the bounded-density condition of \S6.2 again, as it must be: a single superconditioning is a two-outcome update satisfying (M).

\emph{Random update.} The family $\{(\pi_i, Q_i)\}$ is \emph{simultaneously} a conditioning---one extension $(\Omega \times I, P')$ with $P'|_\Omega = P_0$, a partition $\{E_i\}$ with $P'(E_i) = \pi_i$, and $P'(\cdot \mid E_i)|_\Omega = Q_i$---if and only if (M) holds. Necessity is the law of total probability; sufficiency is the construction $P'(A \times \{i\}) := \pi_i\,Q_i(A)$. In the extension, the conditioning event $E_i$ is ``my credence becomes $Q_i$'': the observation is the agent's own new state. (The author attributes this reading to Sam Eisenstat; \ATTR, not independently vetted.) And in the extension, Reflection holds \emph{by construction}: $P'(\cdot \mid E_i) = Q_i$.

\subsection{Two descriptions of one process, with a caveat (\AUTHOR's perspective point; \CLAUDE's caveat)}
The author: the choice between the Bayesian-update representation and the radical-probabilist wider space ``becomes a matter of perspective, two descriptions of one learning process.'' For an agent that has \emph{no} beliefs about its own future credences, this is exactly right: any (M)-satisfying update on $\cE$ is the restriction to $\cE$ of a Bayesian conditioning on $\cE \otimes \sigma(I)$, and the Bayesian on the larger algebra, restricted to $\cE$, is a radical probabilist; superconditioning supplies the missing joint over world and outcome, and the joint it supplies is reflective. But notice that the \emph{independent} joint $P''(A \times \{i\}) := \pi_i P_0(A)$ is also a conservative extension with the same marginal and the same law, and the two extensions give opposite verdicts on whether the update is worth wanting---Good's theorem under one, the pill under the other. So the radical probabilist's state, by itself, does not determine that verdict; superconditioning re-describes the update's \emph{law} for free and its \emph{mechanism} only by a choice. The caveat is the agent that \emph{does} have such beliefs---which is any agent asked whether it wants a modification, since answering requires a joint over the world and the modification's outcome. For that agent the superconditioned joint is faithful only if it coincides with the agent's own joint, which is exactly (R). In the coin example of \S6.3 the agent's joint is the independent one; the superconditioned joint is a different agent's, one for whom the coin is informative. So once the agent has beliefs about how its update is connected to the world, the Bayesian description is not a free re-description but a constraint: it holds iff the agent reflects on its update process. The load therefore sits on the accuracy of those beliefs---the agent's conditional beliefs about the world given each way it might be modified.

\subsection{Other trust principles (\CLAUDE, compact)}
Dorst, Levinstein, Salow, Husic and Fitelson's Total Trust (``Deference Done Better'', 2021, from memory) is the condition $\E_{P_0}[X \mid \E_{P_1} X \ge t] \ge t$ for all bounded $X$ and thresholds $t$, and their theorem is that it is equivalent to preferring, in every decision problem, to act on the source's opinion rather than one's own. Reflection implies Total Trust; Total Trust does not imply (M). Two worlds, $P_0 = (0.3, 0.7)$, source opinion $Q_1 = (0.9, 0.1)$ in world~1 and $Q_2 = (0.1, 0.9)$ in world~2: Total Trust holds (checked on a grid of payoff vectors), but $\E[Q(\{1\})] = 0.34 \ne 0.3$, and the only mixing weights that recover $P_0$ from $\{Q_1, Q_2\}$ are $(\tfrac14, \tfrac34)$, not the agent's $(0.3, 0.7)$. The source here is \emph{modest}---$Q_1$ gives probability $0.1$ to the world in which it would have been $Q_2$---and modesty is where the criteria separate. Consequence: the updates it is rational to want strictly include the superconditionable ones, and the gap is the modest source, which any retraining process is. The body of this note used (R), which is the clean case; which class a theory of legitimate correction should use is left open.

\subsection{The gap in the story so far}
Everything in this section is about $P$ with $U$ fixed; in Jeffrey--Bolker terms it characterizes the desirability changes of the form $V \to \E_{P_1}[U \mid \cdot]$ that Conjecture~5.6 allows. \S2 showed how a change of $U$ becomes a change of $P$ on the extended algebra, and \S3 that on the algebra containing the outcome it was conditioning all along. So the rational value changes are, in this section's terms, the superconditionable updates of the value coordinate---and the gap between (M) and (R) of \S6.3 is exactly the gap between a coin and a teacher.

% ======================================================================
\subsection{Which trust principle, for the framework}
The rule that licenses wanting a modification is reflection with respect to the modification event---(R), (R$^U$), desirability reflection, conditional on legitimacy in the author's framework---or its Total-Trust weakening. No-expected-net-update, (M), is necessary and is what superconditioning needs, but the coin example shows it is not sufficient: an update can satisfy (M) and (Z) and be strictly worse than no update---and be Dutch-bookable, once the outcome is in the agent's algebra. A natural first statement of the legitimacy condition is ``no expected net update with respect to legitimate updates''; this note suggests it should be stated as reflection---a condition on the agent's conditional beliefs given each correction---rather than as the martingale condition on the correction's marginal law, since only the former distinguishes a teacher from a coin.

\section{Open questions}
\begin{enumerate}[label=(\roman*),leftmargin=2.5em]
\item The modest source: Total Trust without (M) (\S6.6) is worth wanting and not superconditionable; whether the framework wants the narrower reflective class or the wider Total-Trust class is unanswered.
\item The richness condition (\S2.2, (Z$^U$)): a value change to a target outside the convex hull of the prior's support is not representable as a rational belief change on any extension; what an agent should do when offered one belongs with changes to the algebra itself, which this note does not treat.
\item The known-target case with a \emph{predictable} source: a capable agent can often predict what the overseers will say, and a predicted correction is a known target; under reflection it has already been adopted, and the residual---corrections predicted and not adopted---are by Claim~1 exactly the ones the agent judges illegitimate. Whether that residual is empty for a well-trained agent is the fully-updated-deference question, which this note locates but does not answer.
\item \S1.4 and Theorem~2.4 assume that the keep/change choice and the acts carry no information of their own about the facts or the source's output. Dropping that---the choice to accept being itself evidence about $\theta$, so that an agent may refuse information it would rather not be seen to want---is the evidential-channel question, and it is untouched here.
\end{enumerate}

% ======================================================================
\appendix
\section{Worlds beyond ultrafilters: designated meets, annihilation, and where utility lives}

\begin{quote}\small
This appendix summarizes (\CLAUDE) a construction from the author's working notes on decision problems (v2.1, revised with Tim Parker; included alongside this note as \emph{decision-problems-v2}), which generalizes the choice of world-set made in \S5.4. Nothing in the body of the note depends on it; it records why the ultrafilter construction is one choice among several and what the alternatives cost.
\end{quote}

\subsection{The construction}
An \emph{event space} is a pair $(\cE, J)$: a Boolean algebra together with a set $J$ of \emph{designated meets}, families $D \subseteq \cE$ whose infimum $\bigwedge D$ exists in $\cE$. Designating $D$ is an optional $\omega$-rule for that family: it asserts that $\bigwedge D$ follows from its conjuncts jointly, whereas an undesignated infimum is a greatest lower bound that is permitted to fail while every conjunct holds. A \emph{world} of $(\cE, J)$ is a set $\omega \subseteq \cE$ satisfying \S5.4's (i)--(iv) and in addition
\begin{enumerate}[label=(\roman*),leftmargin=2.5em,start=5]
\item for every $D \in J$, if every member of $D$ is in $\omega$ then $\bigwedge D \in \omega$.
\end{enumerate}
For finite $D$, (v) is just (iii), so $J$ has only infinitary content, and $J = \emptyset$ gives back exactly the ultrafilters of \S5.4. A \emph{probability} on $(\cE, J)$ is a finitely additive $P$ continuous along designated meets,
\[
P\Bigl(\bigwedge D\Bigr) = \inf\Bigl\{\,P\Bigl(\bigwedge F\Bigr) : F \subseteq D \text{ finite}\Bigr\} \quad\text{for every } D \in J.
\]
On a $\sigma$-algebra of sets with every countable family designated this is exactly countable additivity; with $J = \emptyset$ it is exactly finite additivity. And worlds are exactly the two-valued probabilities on $(\cE, J)$: the displayed condition, read two-valuedly, is (v). So one parameter decides at once which worlds exist and which probabilities count, and the same condition governs both. Designating all countable families, or none, are the by-cardinality special cases.

\subsection{The phenomena}
\emph{Existence at $J = \emptyset$} is \S5.4, steps 2--3: worlds always exist, and finitely additive probabilities upgrade to $\sigma$-additive measures on the Stone space. \emph{Witness-less worlds} are the price: on $2^{\N}$, any ultrafilter extending the cofinite filter contains $\N = \bigcup_n \{n\}$ but no $\{n\}$---some number is selected, and no number is it; in the finite--cofinite algebra the cofinite sets themselves form such an ultrafilter, with no choice principle needed. Whether this is a defect or a resource (an ``unspecified large number'' is exactly what limit constructions want) depends on the subject matter, which is why designation is per family and optional. \emph{Annihilation}: designating too much can leave no worlds at all. On the Lebesgue measure algebra---Borel subsets of $[0,1]$ modulo null sets---with countable meets designated, atomlessness yields a decreasing chain of halves with measures $2^{-n}$ whose infimum in the algebra is $0$; a world following the chain holds every element and by (v) must hold $0$, a contradiction. On $\Borel[0,1]$ itself, by contrast, designating countable meets yields exactly the Dirac worlds $\delta_x$: null events are retained as distinct objects there, and the worlds live on them. The general repair for pointless algebras is Loomis--Sikorski: represent the algebra as a $\sigma$-field of sets modulo a $\sigma$-ideal and read satisfaction modulo negligibility. \emph{The world/probability mismatch}: worlds and probabilities cannot be handed the same maximal designation. A probability additive over all existing joins is purely atomic---$[0,1] = \bigvee_x \{x\}$ exists in $\Borel[0,1]$, and Lebesgue measure would need $1$ to be a supremum of finite sums of $0$---so diffuse belief lives exactly on undesignated joins: pressure pushes worlds toward more designation and states toward less, and where each sits is part of a model's content.

\subsection{Where the Radon--Nikodym utility lives (\CLAUDE)}
Step~5 of \S5.4 returns $U$ as an equivalence class modulo null sets---an element of $L^\infty$ of the measure algebra $\cE / \mathcal N_P$. That object exists whether or not worlds do. At $J = \emptyset$ it can be realized as a function on ultrafilters, witness-less ones included. With countable meets designated on a measure algebra there may be no worlds at all (annihilation) while $U$ still exists, as a function modulo a $\sigma$-ideal: the Loomis--Sikorski object. So the utility random variable natively lives on the \emph{state} side of A.2's mismatch, not the world side; Radon--Nikodym is indifferent to which worlds one admits, and the choice of $J$ is a choice about how to name the values of $U$, not about whether $V$ is the expectation of something. For this note's question that is reassuring: the translation of \S5 does not wait on the foundational choice, and the worked examples are finite, where designation is moot and worlds are atoms.

% ======================================================================
\section*{References (from memory; unverified against text)}
\addcontentsline{toc}{section}{References}
\begin{itemize}[leftmargin=1.5em,itemsep=1pt]
\item Bolker, E. (1966). Functions resembling quotients of measures. \emph{Trans.\ AMS} 124.
\item Bolker, E. (1967). A simultaneous axiomatization of utility and subjective probability. \emph{Philosophy of Science} 34.
\item Jeffrey, R. (1965/1983). \emph{The Logic of Decision}.
\item Jeffrey, R. (1992). \emph{Probability and the Art of Judgment}.
\item Diaconis, P. and Zabell, S. (1982). Updating subjective probability. \emph{JASA} 77.
\item Goldstein, M. (1983). The prevision of a prevision. \emph{JASA} 78.
\item van Fraassen, B. (1984). Belief and the will. \emph{Journal of Philosophy} 81.
\item Skyrms, B. (1987). Dynamic coherence and probability kinematics. \emph{Philosophy of Science} 54.
\item Good, I. J. (1967). On the principle of total evidence. \emph{BJPS} 17.
\item Dorst, K., Levinstein, B., Salow, B., Husic, B. and Fitelson, B. (2021). Deference done better. \emph{Philosophical Perspectives} 35.
\item Broome, J. (1990). Bolker--Jeffrey expected utility theory and axiomatic utilitarianism. \emph{Review of Economic Studies} 57.
\item Joyce, J. (1999). \emph{The Foundations of Causal Decision Theory}, ch.~4.
\item Karni, E. (1985). \emph{Decision Making Under Uncertainty: The Case of State-Dependent Preferences}.
\item Omohundro, S. (2008). The basic AI drives. \emph{Proc.\ AGI 2008}.
\item Bostrom, N. (2012). The superintelligent will. \emph{Minds and Machines} 22.
\item Soares, N., Fallenstein, B., Yudkowsky, E. and Armstrong, S. (2015). Corrigibility. \emph{AAAI Workshop on AI and Ethics}.
\end{itemize}

\end{document}
